\subsection{The remaining nonlocal three-row source of order nine}

\begin{theorem}[Proved: the nine-element retraction source opposite $R_3$]
\label{CP2:A-nine-source8038-all-states}
Let $M$ consist of the identity, a $3\times2$ rectangular band
$K=\{e_{ij}:0\le i\le2,\ 0\le j\le1\}$, and two further
elements $a,b$.  Their products and actors are
\[
 a^2=a,\quad b^2=b,\quad ab=b,\quad ba=a,
 \qquad
 ae_{ij}=be_{ij}=e_{\sigma(i),j},
 \qquad
 e_{ij}a=e_{i0},\quad e_{ij}b=e_{i1},
\]
where $\sigma(0)=\sigma(1)=0$ and $\sigma(2)=2$.
For every finite left-$R_3$, right-$M$ biset, all rational
degree-two products of its collage vanish.  No state-count,
boundedness, or second-Betti-number hypothesis is needed.

This monoid is source $8038$ of the frozen positive order-nine
catalog, with identity label $0$, $e_{ij}=1+2i+j$, $a=7$, and $b=8$.
\end{theorem}
\begin{proof}
Put $R=\mathbb QM$ and set
\[
 e=e_{00},\qquad p=1-a,\qquad
 u=e_{10}-e_{00},\qquad v=e_{20}-e_{00}.
\]
The following is a complete left-projective resolution of the
trivial module:
\[
 0\longrightarrow Re\oplus Ra
 \xrightarrow{d_2}Ra\oplus Rp
 \xrightarrow{d_1}Re\xrightarrow{\epsilon}\mathbb Q
 \longrightarrow0,                                         \tag{1}
\]
where
\[
 d_1(x,y)=xv+yu,\qquad d_2(t,s)=(t,s(b-a)).                 \tag{2}
\]
We verify the entire resolution directly.
The idempotents $e,a,p$ give projective principal left ideals,
with bases
\[
\begin{aligned}
 Re&=\mathbb Q\{e_{00},e_{10},e_{20}\},\\
 Ra&=\mathbb Q\{e_{00},e_{10},e_{20},a\},\\
 Rp&=\mathbb Q\{p,b-a,e_{01}-e_{00},e_{11}-e_{10},
                                     e_{21}-e_{20}\}.
\end{aligned}
\]
Right multiplication by $v$ sends $a$ to $v$ and kills the
three band basis vectors of $Ra$.
Right multiplication by $u$ sends $p$ to $u$ and kills
the other four displayed basis vectors of $Rp$.
Thus the image of $d_1$ is precisely
$\ker\epsilon=\mathbb Qu\oplus\mathbb Qv$.
The first part of $d_2$ is the literal inclusion $Re\subset Ra$.
On its other part, right multiplication by $b-a$ sends
$a$ to $b-a$ and sends $e_{i0}$ to $e_{i1}-e_{i0}$.
It is therefore injective and its image is the four-dimensional
kernel of right multiplication by $u$ on $Rp$.
The two parts of $d_2$ identify its image exactly with
$\ker d_1$.  They also show that every map is left $R$ linear
and has the stated target.  This proves (1), including
injectivity at its left endpoint.

Let $L=\mathbb Q[S]$, retaining its commuting left-$R_3$ action.
Tensoring (1) on the left by $L$ gives the complex
\[
 Le\oplus La\longrightarrow La\oplus Lp\longrightarrow Le.
\]
The first component of the top differential is the inclusion
$Le\hookrightarrow La$, since $ea=e$.
There are no higher resolution terms.  Consequently its second
homology is naturally
\[
 \operatorname{Tor}^{R}_2(L,\mathbb Q)
 =\{s\in La:s(b-a)=0\}
 =La\cap Lb
 =\mathbb QF,\qquad F=Sa\cap Sb.                            \tag{3}
\]
Indeed $sa=s$ for $s\in La$, and $b$ is an idempotent,
so $s(b-a)=0$ means $sb=s$.  The images $La,Lb$ are the
permutation spans of the actual image subsets $Sa,Sb$,
whose intersection is their common permutation span.
Thus $F$ consists exactly of the states fixed by both $a$ and $b$.
It is an actual left-$R_3$ set, because that action commutes
with $a,b$.  The identification (3) is natural for every such
commuting map.

For $L=\mathbb Q$ with the trivial right action, the top
differential in (1) becomes $(t,s)\mapsto(t,0)$.
Hence $H_2(BM;\mathbb Q)\cong\mathbb Q$, represented by the
second top summand, and $H_j(BM;\mathbb Q)=0$ for $j\ge3$.
The actual state augmentation $\mathbb Q[S]\to\mathbb Q$
therefore induces on second source Tor the literal map
\[
 \mathbb QF\longrightarrow\mathbb Q,
 \qquad \sum_{f\in F}c_f[f]\longmapsto\sum_{f\in F}c_f.       \tag{4}
\]

Let $E=B(R_3,S,M)$ be the actual bridge.
The trivial right $\mathbb QR_3$ module has projective dimension
at most one.  The natural bridge spectral sequence is
\[
 E^2_{i,j}=\operatorname{Tor}^{\mathbb QR_3}_i
    (\mathbb Q,\operatorname{Tor}^R_j(L,\mathbb Q))
 \ \Longrightarrow\ H_{i+j}(E;\mathbb Q),
 \qquad 0\le i\le1,\quad0\le j\le2.                         \tag{5}
\]
If $F$ is empty, (3) makes its entire $j=2$ row zero.
Thus $H_3(E;\mathbb Q)=0$.
Both corners have zero fourth and third homology, so the
two-block exact sequence gives $H_4$ of the collage equal
to zero.  This settles that case.

If $F$ is nonempty, the term
\[
 E^2_{0,2}=\mathbb Q\otimes_{\mathbb QR_3}\mathbb QF
\]
is a permanent subquotient in total degree two; in fact it
is the filtration-zero subgroup of $H_2(E;\mathbb Q)$,
because (5) has no columns outside $0,1$.
Naturality identifies its map to $H_2(BM;\mathbb Q)$ with
the ordinary coinvariant map induced by (4).
The coinvariants of the nonempty finite set $F$ identify
basis states along the relations of the action, so their
augmentation onto $\mathbb Q$ is surjective.
It follows that the actual projection
\[
 H_2(E;\mathbb Q)\longrightarrow H_2(BM;\mathbb Q)
\]
is surjective.  Dually, the bridge pullback is injective
on $H^2(BM;\mathbb Q)$.
Since $BR_3$ is contractible, the two-block cohomology
sequence now makes every second class of the collage
relative to both corners.  The actual suspension quotient
kills all products of these classes.

For an empty state set, the collage is the disjoint union
of its corners and (1) gives zero fourth cohomology directly.
The literal table, principal-ideal bases, all maps in (1),
and the natural augmentation in (4) are reconstructed in
\path{converse/scope18/nine_source8038_literal_binding.json}.
\end{proof}
